\section{The category on a nilpotent orbit}\label{orb:section}

Fix an isomorphism $\E\simeq\C$ and a Frobenius-fixed retained
component.  Use the quadratic model of Theorem~\ref{mod:quadratic-model},
with $J=\Aut(\sigma)$, $\mathfrak j=\Lie(J)$ and
$V=H^1(\overline X,\ghat_\sigma)$.  All representations in this section
are algebraic; rational representations of an algebraic group may be
unions of finite-dimensional representations.  The group $J$ is reductive
but need not be connected.  We first omit the completion along the
invariant null fiber.  Section~\ref{tr:section} restores that condition,
on both traces and their actual compact representatives.

On each nilpotent orbit we construct compact generators, extend them
with Weil structures to its closure, and compute their morphisms. With
pure coefficient normalizations, the generators and their weight-zero
morphism cohomology form a semisimple category, and the remaining weights for the
inverse action $S$ are strictly negative. Section~\ref{tr:section} uses
these facts to contract the nonconstant part of each orbit trace.

\subsection{The curved model and supported localization}

Put $R=\C[V]$.  The derived zero fiber of the moment map is represented
by the Koszul differential graded algebra
\begin{equation}\label{orb:derived-koszul}
 A=\bigl(R\otimes\Lambda(\mathfrak j[1]),d_\mu\bigr),
 \qquad |\eta_a|=-1,\qquad d_\mu\eta_a=\mu_a.
\end{equation}
Here a basis of $\mathfrak j$ indexes the components $\mu_a$ of
$\mu:V\to\mathfrak j^*$.  In particular, no regularity of this section
is asserted or needed.  Its graded Koszul dual is the curved algebra
\begin{equation}\label{orb:potential}
 B=\C[V\times\mathfrak j],\qquad
 |x^*|=0,\quad |z^*|=2,\qquad
 W(x,z)=\langle z,\mu(x)\rangle.
\end{equation}
A curved module has a differential $\delta$ of degree one satisfying
$\delta^2=W$.  A simultaneous reversal of the curvature sign gives the
same calculations after the corresponding convention change.

\begin{lemma}[Koszul model and support]\label{orb:koszul}
The compact ind-coherent category on the uncompleted derived quotient
$[\mu^{-1}(0)/J]$ is the category of finite coherent graded
$J$-equivariant factorizations of \eqref{orb:potential}.
Coherent singular support in a closed cone $\Omega\subset\mathfrak j$
corresponds to local vanishing of the factorization off
$V\times\Omega$.  The completion in
Theorem~\ref{mod:quadratic-model} additionally imposes ordinary support
on the inverse image of $0\in\Spec R^J$.
\end{lemma}

\begin{proof}
The uncompleted derived quotient is a QCA stack: it is
quasi-compact, with affine stabilizers, and its classical inertia is of
finite presentation, as for every finite-type quotient here. Its compact
ind-coherent objects are therefore its coherent objects by
\cite[Theorem~3.3.5]{DGFiniteness}.  We recall the Koszul construction
to specify the derived and grading conventions.
A coherent $A$-module is perfect as an $R$-complex because
$R$ is regular.  Choose a bounded finite projective $R$-resolution and
transfer the actions of the degree-minus-one generators to it, allowing
higher homotopies.  The Koszul twisting cochain, on the symmetric
\emph{divided-power coalgebra} of the dual generators, turns those
homotopies into a differential on its extension to $B$.  Its square is
$\sum_a z_a\mu_a$.  The linear coefficients encode nullhomotopies of
$\mu_a$, and the higher coefficients encode all compatibility
homotopies for the exterior-generator actions.  Conversely, the
coefficients of such a curved differential, specialized at $z=0$,
recover that homotopy-coherent $A$-action.

The usual free Koszul resolution proves that these constructions are
inverse on derived modules and morphisms: after filtering by exterior
length, it is the ordinary exterior/symmetric Koszul equivalence, and
the equation differential contributes precisely the curvature term.
There is no completed polynomial ring here.  Between bounded finite
projective models only finitely many $z$-monomials can occur in any
fixed total degree.  Derived morphisms on the curved side are sections
of differential Hom; its two curvature terms cancel.

One may equivalently define coherent curved modules modulo
totalizations of finite exact sequences.  On the smooth ambient
space, finite locally free replacements compute this category.
They are obtained by curved free covers, with the partner part of the
differential correcting the square to $W$; the underlying coherent
resolutions terminate by regularity.  Reductivity permits equivariant
choices for $J^\circ$, followed by equivariance for the finite component
group.  This is also the equivariant derived-zero-locus Koszul
equivalence of \cite[Theorem~2.3.3]{Toda}, with its differential graded
formulation in \cite[Theorem~2.1]{PT-koszul}.

The variables $z_a$ are the degree-two complete-intersection operators
on the $A$-side.  Their localization is the localization defining
coherent singular support.  Under the equivalence, the category with
that singular support is exactly the kernel of restriction to the
complementary open, with the corresponding quotient description;
this is \cite[Proposition~2.3.9]{Toda}.  Support here is support of the
object in the curved category: a locally free underlying module can
represent an object vanishing on an open set.

The remaining completion imposes a different, ordinary support
condition.  For a DG scheme almost of finite type,
\cite[Proposition~7.4.5]{DGIndschemes} identifies ind-coherent sheaves
on its formal prestack completion with the kernel of restriction to
the complementary open; its compact objects are the coherent objects
with that ordinary support by \cite[Proposition~4.1.7(b)]{IndCoh}.
Smooth descent, \cite[Corollary~10.4.5]{IndCoh}, identifies the
presentable supported categories on the quotient by $J$. Compactness
on the quotient also uses the QCA theorem just invoked and the continuous
support right adjoint of \cite[Corollary~4.1.5]{IndCoh}. If $i$ denotes
the supported inclusion and $\Gamma$ its continuous right adjoint, then
$\Hom(iM,-)=\Hom(M,\Gamma(-))$ shows that $i$ preserves compactness.
Conversely, since $i$ preserves colimits, an ambient compact whose
restriction to the complementary open vanishes is compact in the
supported category. Its compact objects are therefore precisely the
supported coherent objects. This identifies compacts after presentable
descent, without commuting compact objects with a descent limit.
Applied to the invariant null fiber in our derived affine model, it
proves the last condition in the lemma.
The completion is the formal prestack, not the spectrum of a completed
coordinate ring.  This ordinary support condition leaves the
singular-direction operators unchanged and can be intersected with
the prescribed $z$-support.
\end{proof}

Write $\mathscr B_\Omega$ for the resulting compact category with
$z$-support in $\Omega$, and $\mathscr B_{\Omega,0}$ when the invariant
null-fiber condition is also imposed.  We use the same notation with
$\operatorname{Ind}$ when a presentable category is required for a trace.  All
compact quotients below are idempotent completed.

\begin{lemma}[Supported localization]\label{orb:localization}
Let $\Omega'\subset\Omega$ be closed invariant cones.  Restriction to
the complementary open gives the compactly generated localization
with compact quotient $\mathscr B_\Omega/\mathscr B_{\Omega'}$.
If $\mathcal O$ is an orbit closed in that open, its supported compact
category is generated by coherent curved pushforwards from
$V\times\mathcal O$.  Morphisms are computed in the open before taking
derived equivariant sections.  The statements also hold with invariant
null-fiber support.
\end{lemma}

\begin{proof}
Choose finite homogeneous equivariant generators of the relevant
ideal sheaves on a finite affine cover.  Koszul complexes on their
powers, and their duals, give the ordinary local-cohomology projector
and complementary localization.  Tensoring a curved module by these
ordinary complexes is defined because their differentials have square
zero.  A compact object's identity factors through a finite stage of
the projector if the object is supported on its center.  Thus it is a
summand of a finite supported stage.

Let $i:Z_0\hookrightarrow U$ be the smooth closed locus in the
smooth ambient open, with ideal sheaf $I$, and let $K_n$ be the finite
ordinary perfect projector stage just chosen.  Its locally free terms
need not be supported on $Z_0$.  Its bounded ordinary cohomology
sheaves are coherent, supported on $Z_0$, and annihilated by powers
of $I$.  Since $U$ is smooth, those coherent sheaves are perfect.
Ordinary truncation triangles therefore express $K_n$ in their thick
span.  Each such sheaf has a finite $I$-adic filtration whose
successive quotients have the form $i_*F$, for coherent sheaves $F$
on $Z_0$.  The truncations, ideal filtrations, and their maps retain
equivariance and grading.

Ordinary perfect complexes act exactly by tensor product on coherent
curved modules.  The extra ordinary differential has square zero,
so this action takes the preceding finite triangles to curved
triangles without changing the potential.  For a locally free curved
representative $M$, the projection formula gives
\[
 (i_*F)\otimes M\simeq i_*\bigl(F\otimes\mathbf L i^*M\bigr).
\]
The right side is a coherent curved pushforward with the restricted
potential.  Consequently $K_n\otimes M$ lies in the thick span of
such pushforwards.  The factorization of the identity through this
stage makes $M$ a summand, proving generation.  The ordinary
truncations have been applied only to $K_n$, before its exact action
on the curved factor.

The same local-cohomology construction computes restriction and the
quotient on morphisms; the support/quotient compatibility is also the
one in \cite[Proposition~2.3.9]{Toda}.  On the ambient open, finite
covers and finite \v{C}ech descent compute the ordinary perfect-complex
statements.  On a homogeneous orbit, equivariant descent takes the
finite Lie-cochain complex for its unipotent stabilizer, followed by
exact invariants for the reductive quotient.  These bounded operations
commute with the finite triangles and direct-sum folding of the
grading.  The argument also applies on homogeneous vector bundles.
Ordinary support in the invariant base is preserved at every step.
\end{proof}

\subsection{Gradings and Clifford generators}

The intersection of $\mathfrak j$ with $\Nhat$ is its nilpotent cone:
for a reductive subgroup, the semisimple and nilpotent parts of an
element agree with those in the ambient faithful representation.
There are finitely many nilpotent $J$-orbits.  Fix one, $\mathcal O=J.e$,
and choose an $\mathfrak{sl}_2$-triple $(e,h,f)$ in $\mathfrak j$.
Set
\begin{equation}\label{orb:notation}
 \lambda=\lambda_e,\qquad
 C_e=J^e=L_e\ltimes U_e,\qquad
 L_e=Z_J(\mathrm{SL}_{2,e}),\qquad
 N_e=\mathfrak j/[\mathfrak j,e].
\end{equation}
The zero orbit is included, with trivial triple and $U_e=1$.
Subscripts $i$ denote $h$-weights, independently of cohomological
degree or Frobenius weight.  The Lie algebra of $U_e$ has strictly
positive $h$-weights.  The quotient $N_e$ consists of lowest-weight
vectors and has $h$-weights at most zero.

Suppose first that the orbit is Frobenius-stable.  By
Lemma~\ref{mod:compact-normalization}, the weight-zero normalized
transporter can be chosen to fix the triple, preserve a compatible
compact form, and remain of absolute weight zero after compact
$L_e$-twists.  Multiplication by $\lambda(\sqrt q)$ corrects the scalar
motion of $e$ and gives an action fixing $e$.  Write $S_e$ for this
translated inverse action on the Hom calculation.  For an orbit or a
finite list of labels stable only under a power, the same convention
is made for that power.  Absolute weights are measured in units
$q^{1/2}$ (or $q^{m/2}$ for the $m$th power).

There is a separate translation of the grading.  The original grading
acts on $z$-points by $t^{-2}$; lifting it by $\lambda(t)$ makes $e$
fixed.  We call the resulting degree on a fiber at $e$ the
\emph{lifted degree}.  The data needed below are
\begin{equation}\label{orb:weight-table}
\begin{array}{c|c|c}
 \text{factor}&\text{lifted degree before an Ext shift}
                  &\text{absolute }S_e\text{-weight}\\ \hline
 (V_i)^*\text{, as an }x\text{-function}&-i&-1-i\\
 (N_e)_i\text{, as a normal tangent}&i-2&-2+i\\
 \Lie(U_e)_i^*&-i&-i.
\end{array}
\end{equation}
Translating by a group element changes none of the final equivariant
invariants.  In particular these are the actual weights there, rather
than weights for a newly chosen categorical Frobenius.

The original super sign must be retained during this change of
grading.  It differs from the lifted-degree sign by the action of
$\lambda(-1)$.  Consequently the lifted total-degree sign computes
true alternating traces \emph{after taking invariants}.  In
particular, a polynomial coordinate of odd lifted degree remains a
commuting coordinate; it is not replaced by an exterior generator.

Put $Q_e(x)=W(x,e)$.  With the symplectic pairing on $V$, this is, up
to the fixed moment-map sign, $\frac12\langle ex,x\rangle$.
Its restriction to $V_{-1}$ is nondegenerate: $e:V_{-1}\to V_1$ is an
isomorphism and the symplectic pairing identifies $V_1$ with
$V_{-1}^*$.  Let
\[
 D_e=\operatorname{Cl}(V_{-1},Q_e)
\]
be the ordinary Clifford algebra, with $c(v)^2=Q_e(v)$, and let
$\mathsf E_e$ be the category of finite-dimensional $D_e$-modules with
compatible algebraic $L_e$-action.

The related spectral Whittaker construction uses Clifford modules for a
quadratic cohomology complex \cite[Sections~4.5.2--4.5.5]{Raskin};
its orbit trace comparison is Theorem~J in Section~4.6.2 there.
Here we keep explicit generators and their Hom weights, as well as
compact Weil extensions, for the later bilinear tests.

For $P\in\mathsf E_e$, give its coefficient space lifted degree zero.
Its original coefficient parity is the involution
\begin{equation}\label{orb:parity}
 \gamma_P=\rho_P(\lambda(-1)),\qquad
 \gamma_Pc(v)\gamma_P^{-1}=-c(v)\quad(v\in V_{-1}).
\end{equation}
Define $M_P$ over $e$ as the curved pushforward from
$V_{\ge-1}\subset V$ of
\begin{equation}\label{orb:generator}
 \bigl(\mathcal O_{V_{\ge-1}}\otimes P;\quad
                  \delta=c(x_{-1})\bigr).
\end{equation}
The unipotent radical acts on $P$ through its trivial quotient, and
acts trivially on the projection of $V_{\ge-1}$ to $V_{-1}$.
Only quadratic terms with both inputs of weight minus one contribute
to $Q_e$ on this subspace.  Thus \eqref{orb:generator} has square $Q_e$, is
$C_e$-equivariant, and has lifted differential degree one.
The lifted grading and $C_e$-equivariance descend to the required
original graded object over $J.e$.  Denote its curved pushforward into
a surrounding open by the same symbol $M_P$.

\begin{proposition}[Orbit generation]\label{orb:generation}
The objects $M_P$, for simple $P\in\mathsf E_e$, generate the compact
category supported on $J.e$ in an open where it is closed.  Here
``generate'' allows finite cones, shifts and summands.
The category $\mathsf E_e$ is semisimple.
\end{proposition}

\begin{proof}
For a locally free factorization, differentiation of $\delta^2=Q_e$
gives
\[
 \delta(\partial_v\delta)+(\partial_v\delta)\delta
                         =\partial_vQ_e.
\]
Thus the derivatives of $Q_e$ annihilate the identity up to homotopy.
Their common zero locus is $\ker(e:V\to V)$, which consists of
highest-weight vectors and lies in $V_{\ge0}$.  Lemma~\ref{orb:localization}
therefore reduces generation to curved pushforwards from $V_{\ge0}$,
where the potential is zero.

On this affine space, finite graded free modules with finite
stabilizer-representation labels and their degree shifts generate the
zero-potential category.  To justify this with the nonreductive
stabilizer, note first that all finite rational $C_e$-modules have a
finite filtration with $U_e$-trivial subquotients: take unipotent
invariants and iterate, retaining $L_e$-stability.  Rational
$U_e$-cohomology has cohomological dimension $\dim U_e$ and is computed
by its finite Lie-cochain complex.  The usual equivariant free tests
therefore detect underlying graded cohomology: tensor by all finite
representation labels, pass to their union in the regular
representation, and then take invariants.  These operations preserve
colimits.  The tests are compact and hence generate.

There is no additional curved object at zero potential invisible to
this test.  After a finite projective replacement, an acyclic folded
differential module over the regular polynomial ring is absolutely
acyclic: resolve its cycles and boundaries by finite projective
resolutions and fold the resulting bounded exact complexes.
Equivariant descent adds only the bounded Lie-cochain complex above.
Equivalently, the lifted coordinate degrees on $V_{\ge0}$ are
nonpositive; filtering finite free representatives by generator degree
reduces to finite complexes over the degree-zero polynomial ring.

Stabilizing the zero section of the nondegenerate quadratic space
$V_{-1}$ replaces a generator on $V_{\ge0}$ by
\eqref{orb:generator} with the regular Clifford module.  In coordinates
this is the Koszul differential for $x_{-1}=0$, together with its
Clifford partner; their square is $Q_e(x_{-1})$.  The normal coordinates
have lifted degree one, so the Koszul shifts put the coefficient
module in lifted degree zero.  Tensoring by $L_e$-representations and
allowing shifts gives all the preceding generators.

Finally $D_e$ is a finite-dimensional semisimple complex algebra,
including when $\dim V_{-1}$ is odd.  A $D_e$-linear splitting of any
short exact sequence in $\mathsf E_e$ can be averaged over a compact
form of the reductive group $L_e$.  Conjugation preserves $D_e$-linearity,
so the averaged splitting is also $L_e$-equivariant.  This proves
semisimplicity, including the finite component group, and reduces the
generating list to its simple objects.
\end{proof}

\subsection{Actual extensions with Weil structures}

\begin{proposition}[Compact orbit-closure lifts]\label{orb:weil-lifts}
Every $M_P$ is the restriction of a compact object with $z$-support in
$\overline{J.e}$.  If $P$ has a compatible locally finite cyclic
structure, the extension can be given that structure.  For a simple
label fixed by Frobenius, a structure may be chosen whose
triple-translated coefficient eigenvalues have modulus one.
The same assertions hold for a power of Frobenius, for a finite cyclic
list of labels, and after tensoring $P$ by any finite-dimensional
$L_e$-representation with compatible cyclic structure, placed in
lifted degree zero.  Inverse structures are available for opposite
category classes.
\end{proposition}

\begin{proof}
Let $P_e=J_{\ge0}$ be the subgroup of elements for which
$\lim_{t\to0}\lambda(t)g\lambda(t)^{-1}$ exists, and put
$J_0=Z_J(\lambda)$.  The Jacobson--Morozov incidence map is
\begin{equation}\label{orb:incidence}
 \pi:J\times^{P_e}\mathfrak j_{\ge2}
                       \longrightarrow\overline{J.e},
 \qquad [g,z]\longmapsto\operatorname{Ad}(g)z.
\end{equation}
Its connected version is the resolution in
\cite[Section~2.3]{Namikawa}; the following argument records the
properties, including disconnected groups, that we use.
The $\mathfrak{sl}_2$ decomposition gives
$[\mathfrak j_{\ge0},e]=\mathfrak j_{\ge2}$.  Hence
$P_e^\circ.e$ is open dense in $\mathfrak j_{\ge2}$.  Every component
of $P_e$ preserves this unique open $P_e^\circ$-orbit, so
$P_e.e=P_e^\circ.e$.  The centralizer decomposition in
\eqref{orb:notation} gives $J^e\subset P_e$.
Consequently the open part of the source of \eqref{orb:incidence} is
$J\times^{P_e}(P_e.e)=J/J^e$, mapped isomorphically to $J.e$.
The entire source is equidimensional of dimension $\dim(J/J^e)$.
Its complementary closed subset has smaller dimension and has
$J$-invariant image, so that image cannot meet $J.e$.
Thus this is the whole inverse image of the orbit.
Finally $J/P_e$ is a finite disjoint union of flag varieties; the
incidence space is closed in $(J/P_e)\times\mathfrak j$.
The map is therefore proper and its image is $\overline{J.e}$.

The degree-two component of this construction uses the open orbit
$J_0.e\subset\mathfrak j_2$, whose stabilizer is $L_e$.
Over that open orbit the module $P$ gives an equivariant coherent module
for the family of Clifford algebras of
\[
 Q_z|_{V_{-1}},\qquad z\in\mathfrak j_2.
\]
This Clifford algebra is finite as a module over
$\mathcal O_{\mathfrak j_2}$, also where the quadratic form degenerates.
Extend the module coherently over $\mathfrak j_2$ as follows.  Its
quasi-coherent direct image from the open orbit is a rational
$J_0$-module.  Choose finitely many module generators over that open,
enlarge them to a finite-dimensional equivariant space, and take their
span under the finite Clifford algebra and the polynomial ring.
This is a coherent equivariant submodule whose restriction is the
original family.  More explicitly, let $t$ act on this coefficient
family through the central element $\lambda(t^{-1})\in J_0$.  It sends
$z\mapsto t^{-2}z$ and
$c_z(w)\mapsto c_{t^{-2}z}(tw)$.  With the original $x$-points fixed,
$c_z(x_{-1})$ consequently has degree one.  At $t=-1$ its coefficient
parity is precisely $\lambda(-1)$.  Thus equivariant extension also
preserves the required original grading.

Inflate this extension along
$\mathfrak j_{\ge2}\to\mathfrak j_2$ and tensor it with
$\mathcal O_{V_{\ge-1}}$.  The differential $c(x_{-1})$ has square
$W(x,z)$ on this locus: a term pairing $z_k$ with $x_i,x_j$ requires
$k+i+j=0$, with $k\ge2$ and $i,j\ge-1$, and hence
$k=2$, $i=j=-1$.  The construction is $P_e$-equivariant because its
unipotent radical acts trivially on the two indicated associated
graded quotients.  Induction to $J$ and proper pushforward from
\[
 J\times^{P_e}(V_{\ge-1}\times\mathfrak j_{\ge2})
                 \longrightarrow V\times\mathfrak j
\]
give the required coherent factorization.  This map is proper by the
same closed-incidence argument.  Over $J.e$ its restriction is exactly
$M_P$, so this is an actual lift, with the asserted support.

For completeness, the equivariance in this extension can include a
chosen cyclic operator, rather than merely its action on isomorphism
classes.  Take the Zariski closure of the group generated by the
algebraic equivariance data, the grading, and that operator; if
necessary take a common algebraic cover acting on the coefficient
module.  Rational direct image is the union of finite-dimensional
representations of this group.  Thus the finite generating space in
the preceding paragraph can be chosen stable under the operator, and
the entire extension and proper pushforward retain its structure.

Here is also why the pure coefficient choices exist.  Use the
finite-type normalizer of the linear and framed data in
Lemma~\ref{mod:compact-normalization}; the action is an algebraic action
of this finite-type group, not of an abstract automorphism group of a
torus.  For a fixed simple equivariant Clifford module, pairs
consisting of a data automorphism and an implementing linear map form
an algebraic group over the stabilizer of its isomorphism class.  Its
kernel is the scalars, by Schur's Lemma.  Given a lift $(u,A)$, take its
semisimple Jordan part and the diagonalizable algebraic group it
generates.  Characterwise polar decomposition in this group separates
its positive-real and unit-modulus parts.  The positive-real part
projects to the identity: the eigenvalues of $u$ on the faithful
normalized data realization have modulus one.  The scalar-kernel
assertion therefore forces this positive-real part to be $(1,cI)$
for some $c>0$.  All eigenvalues of $A$ have the same modulus $c$,
and $c^{-1}A$ is the required pure implementing map.  The Jordan
unipotent part supplies the compatible commuting unipotent lift.

The label orbits are finite as well.  Encode compatible Clifford and
$L_e$ actions as representations of the reductive group
$D_e^\times\rtimes L_e$.  A simple compatible module is irreducible
for this group, because the units linearly span $D_e$.  A connected
algebraic family of automorphisms preserves each irreducible class:
on a fixed maximal compact group, its multiplicities against
irreducible characters are continuous integer-valued functions of the
family, hence constant.  The compact group is Zariski dense, so this
identifies the algebraic representations.  Only the finite component
group of the data normalizer can therefore permute labels.  The
preceding choices apply to a power, or around a finite orbit; a label
fixed by Frobenius needs no power.  Undoing the triple translation
gives the original locally finite Weil structure.  Tensor products
and inversion of the implementing maps preserve the extension
construction, proving the remaining assertions.
\end{proof}

\subsection{Morphisms, parity, and contracting weights}

\begin{proposition}[The orbit Hom calculation]\label{orb:hom-calculation}
For $P,P'\in\mathsf E_e$, the Hom complex between the orbit generators
has a finite equivariant calculation with Euler terms
\begin{equation}\label{orb:hom-euler}
 \left[
 \C[\ker(e:V\to V)]\otimes
 \Hom_{D_e}(P,P')\otimes
 \Lambda^\bullet(N_e[-1])\otimes
 \Lambda^\bullet(\Lie(U_e)^*[-1])
 \right]^{L_e}.
\end{equation}
Here the coefficient factor has lifted degree zero; the shifts in the
normal and Lie factors add one to the degrees in
\eqref{orb:weight-table}.  The formula is an equality of alternating
equivariant characters, and a finite-filtration calculation bounding
weights, rather than an assertion that the Hom complex splits as the
displayed tensor product.

For pure coefficient normalizations, all nonconstant terms have
strictly negative absolute $S_e$-weight.  The weight-zero cohomology is
\begin{equation}\label{orb:weight-zero-hom}
 H^\bullet\Hom(M_P,M_{P'})_{\mathrm{wt}=0}
       =\Hom_{D_e,L_e}(P,P')[0],
\end{equation}
with ordinary composition.  Before the finite invariant calculation,
weight pieces are finite-dimensional and their dimensions grow at most
polynomially over bounded-length weight intervals.  In particular,
the character sums in \eqref{orb:hom-euler} are absolutely convergent
on every normalized compact $L_e$-coset, uniformly on that coset.
\end{proposition}

\begin{proof}
Take an $L_e$-stable lowest-weight slice transverse to $J.e$ in the
$z$-space.  Its tangent is $N_e$; moving this slice by $J$ is smooth
near $e$, so it computes the equivariant transverse fiber of Hom.
Then take $U_e$-cochains and finally exact $L_e$-invariants.

The second input is pushed forward from the equations $n=0$ in the
slice and $x_{<-1}=0$.  A finite locally free stabilization of the
first input is obtained by tensoring its Clifford differential with
the Koszul factorizations of these equations and their potential
partners.  Indeed the difference between the full potential and
$Q_e(x_{-1})$ is a sum of an equation times its partner.  The terms in
$n$ have quadratic partners in $x$; the terms $x_i$ for $i<-1$ have
partners of grade $-i-2\ge0$.  This explicit correction makes the
square of the differential equal to the full potential.
Hom into the second pushforward sets the normal equations to zero.
The $z$-normal Koszul generators remain as exterior \emph{normal
tangent} factors $\Lambda(N_e[-1])$.

Filter finitely by those exterior generators.  The remaining
$x$-Koszul differential pairs every grade $i<-1$ with its grade
$-i-2$ partner.  On an irreducible $\mathfrak{sl}_2$ summand these
pairings are nondegenerate except for its highest-weight vector.
Thus their cohomology is precisely $\C[\ker e]$.  The terms involving
$n$ may give further differentials in this finite filtration, but do
not change its Euler terms.

It remains to calculate the nondegenerate quadratic part.  On
$\C[V_{-1}]\otimes\Hom(P,P')$ its differential is the sum of the
coordinate functions times the Clifford supercommutators.  Choose an
orthonormal Clifford basis, and on a homogeneous coefficient map $a$
put
\[
 L_i(a)=c_i'a,\qquad R_i(a)=(-1)^{|a|}ac_i,\qquad
 T_i=L_i-R_i,\qquad K_i=\tfrac14(L_i+R_i).
\]
Clifford anticommutation gives
\[
 \{T_i,T_j\}=\{K_i,K_j\}=0,\qquad
 \{T_i,K_j\}=\delta_{ij}.
\]
Thus the $K_i$ are contracting partners for the exterior operators
$T_i$.  The differential $\sum_i x_iT_i$ is the polynomial Koszul
calculation: its cohomology is the constant common kernel of the
$T_i$, namely the Clifford super-intertwiners, in lifted degree zero.
All positive polynomial degrees are acyclic.  This calculation works
in even and odd Clifford dimensions alike.

We explain its parity before taking invariants.  A homogeneous map
$a:P\to P'$ of original parity $p$ is a super-intertwiner precisely
when
\[
 a c(v)=(-1)^p c'(v)a.
\]
The map $a\mapsto a\gamma_P^p$ identifies that vector space with
ordinary Clifford intertwiners.  It respects the $L_e$-action and the
triple-preserving cyclic action, because $\lambda(-1)$ is central in
$L_e$ and is fixed by that action.  Only a representation-space
identification is claimed before taking invariants.  An $L_e$-invariant
map commutes with $\gamma_P,\gamma_{P'}$, so an invariant
super-intertwiner has even parity.  After invariants the identification
therefore has its ordinary composition and is exactly
$\Hom_{D_e,L_e}(P,P')$.  This proves that no odd Clifford morphism is
being silently retained in degree zero.

Finally rational $U_e$-cohomology is computed by its finite
Chevalley--Eilenberg complex, adding the terms
$\Lambda(\Lie(U_e)^*[-1])$, with their Lie differential and action on
the other terms.  This proves \eqref{orb:hom-euler} as an alternating
character identity with a finite cohomological filtration.

The weight assertions now follow factor by factor.  Functions on
$\ker e$ have weights $-1-i\le-1$ because its $h$-weights satisfy
$i\ge0$.  The normal factors have weights $-2+i\le-2$ because their
$h$-weights satisfy $i\le0$.  Nontrivial Lie-cochain factors have
weights $-i<0$.  The pure constant coefficient Hom has weight zero.
Thus nothing of any other degree has weight zero, and no differential
can enter or leave its invariant subspace.  This proves
\eqref{orb:weight-zero-hom}.

There are finitely many polynomial generators, all with strictly
negative weights, and finitely many exterior and coefficient factors.
Their monomial counts in bounded-length weight intervals have
polynomial growth.  Multiplication by the geometric decay of the
absolute eigenvalues makes their traces summable.  The same estimates
hold uniformly with any compact $L_e$-twist by
Lemma~\ref{mod:compact-normalization}.  Unipotent parts may be retained:
trace depends on eigenvalues with multiplicity, not on an assertion
that those operators are unitary.
\end{proof}
